% chapters/chapter06.tex
\chapter{Le regole del convitto}
«Muori Hanzaki!»\\
Giovanni schiva per un centimetro un fendente che finisce la sua corsa contro la scala.
«Sei lento, sei sempre stato lento... Non mi serve vederti, posso anticipare ogni tua mossa: questo è il segreto di Kage.»
\par\medskip
«Allora ferma questo colpo!» grida, e scatta brandendo il bastone sopra la testa.
\par\medskip
I due cani lo guardano avvicinarsi e raggiunge la giusta distanza, «Mi hai lasciato raggiungere \textit{ma-ai}, ora non potrai fermarmi.» I suoi movimenti si fanno prima lenti e poi  veloci, ma quando il bastone ha ormai raggiunto  Giovanni lui abbassa di poco la testa, alza le mani e rapidamente lo blocca.
\par\medskip
Osservandolo sembra che stia pregando. Ha le mani chiuse, palmo contro palmo, e in mezzo c'è il bastone, questo è un vera  contromossa ninja.
\par\medskip
\begin{center}
  \includegraphics[width=0.9\linewidth]{media/ma ai senza sfondo.png}
\end{center}

\par\medskip
Ipparchia si avvicina alla scena, lo sconosciuto depone il bastone e le accarezza il pelo sotto la gola. Giovanni accarezza invece la sua testa, scompigliandogli i capelli:
«Dai, adesso aiutami. Prendi il pacco, Luca. Hai visto come ho fatto tardi...»\\
«Infatti, ma questo cane? Dove lo hai preso?»\\
«Non lo so. Ha seguito Ipparchia.»\\
«E come si chiama?»\\
«Non lo so.»

\par\medskip
Luca si ferma e controlla:
«Ecco, è scritto sul collare, Rocky.»

\par\medskip
Rocky lo guarda.\\
«C’è un numero da chiamare?»\\
«Sì, è qui, guarda.»

\par\medskip
Passa qualche istante.\\
«Scusa, comunque c’è.»\\
«Bene, allora domani lo accompagneremo dai vigili. Ci penseranno loro. Portiamolo dentro per ora.»\\
«Lo leghiamo?»

\par\medskip
Giovanni attende. Non risponde subito.\\
«No, prima la libertà.»

\par\medskip
È un grande onore quello che Giovanni ha riservato a Luca: portare il pacco dentro al convitto. Certo, è Giovanni che ha fatto la strada maggiore, ma portarlo dentro significa qualcosa. Gli occhi nel buio lo scruteranno e sapranno che lui, Luca, ha portato il mangiare dentro.

\par\medskip
Però la strada è ancora in salita. Percorrere quella vecchia scala ferrata aggrappandosi con entrambe le mani è un conto; reggere in equilibrio quel pacco di viveri, cercando di non volare giù, è un altro. Ogni passo può tradirlo. La ruggine ha corroso i gradini e, un po’ per il buio, un po’ per il pacco che gli impedisce la vista, Luca deve andare a tentoni per non infilare il piede in una trappola.

\par\medskip
Giovanni, a questo, è abituato.

\par\medskip
«Cosa hai preso di buono, Giovanni?» gli domanda.

\par\medskip
«Una settimana di vita.»

\par\medskip
«Stai attento a non inciampare.»

\par\medskip
Quelle parole suonano quasi come una profezia.

\par\medskip
Un rumore secco, forse una sbarra che cede. Luca sta per perdere l’equilibrio.

\par\medskip
Giovanni si gira di scatto e lo afferra con precisione, evitandogli la caduta.

\par\medskip
«Wow!» esclama Luca. «Grazie!»

\par\medskip
«Tranquillo. Erano solo pochi metri. Però è sempre meglio evitare di cadere.»

\par\medskip
«Già.»

\par\medskip
«Eccoci.»

\par\medskip
Giovanni allunga una mano e trova a tastoni il telaio della porta finestra.

\par\medskip
«Siamo arrivati al primo piano. Aspetta...» Giovanni afferra qualcosa che gli sbarra la strada. Lo tira leggermente, poi ne tasta il contenuto: un carrello della spesa carico di stracci e prodotti per le pulizie blocca l’entrata.
Non è un carrello in buone condizioni. Mancano diverse sfere dai cuscinetti delle ruote, ma sarebbe comodo per portare la spesa. Ormai di carrelli non se ne trovano più. Giovanni però non ha mai pensato di usarlo a quello scopo.

\par\medskip
«Meglio salire al secondo piano» sentenzia.\\
«Ci sono, le pulizie?»\\
«Mi sa di sì.»\\
«Ok, saliamo.»\\
«Ci segue anche Rocky?»\\
«Sì, è dietro Ippa.»\\
«Speriamo non facciano storie… già si lamentano di lei.»

\par\medskip
«Dovrebbero essere in aula uno. Una volta entrati scendiamo dritti per lo scalone, non lo vedranno neanche.»\\
«Hai ragione.»

\par\medskip
Giovanni rallenta il passo per assicurarsi del terreno sotto i piedi, perché non conosce bene la seconda rampa delle scale. Entrare al primo piano è molto più semplice, ma adesso le casalinghe stanno facendo le pulizie, e a memoria di convitto, nessuno ha mai calpestato dove le casalinghe stanno pulendo.

\par\medskip
Sarebbe un errore irreparabile, perché significativo di ignoranza delle regole elementari di appartenenza a una comunità rispettabile. Basta riconoscere alcuni segnali elementari per non creare situazioni problematiche, e uno di questi è la presenza del carrello dei prodotti: un segnale universale di pulizia in corso. Impensabile spostarlo per passare.
